% !TeX encoding = UTF-8
% Formatting-only conversion; see reports for source provenance.
\SourceChapter{2}{فصل دوم: مبانی نظری و پیشینه پژوهش}
% SOURCE Introduction_Chapter2.txt:1-1 [chapter-title-in-wrapper]


% SOURCE Introduction_Chapter2.txt:2-2 [spacing]


% SOURCE Introduction_Chapter2.txt:3-3 [heading]
\SourceHeading{1}{۲{-}۱. مقدمه}

% SOURCE Introduction_Chapter2.txt:4-4 [spacing]


% SOURCE Introduction_Chapter2.txt:5-5 [paragraph]
برای طراحی و پیاده‌سازی یک سیستم امنیتی جامع که محرمانگی تصویر تشخیصی پزشکی و صیانت از فراداده‌های هویتی بیمار را به صورت هم‌زمان تضمین نماید، شناخت عمیق مبانی نظری و تحلیل نقادانه پیشینه پژوهش امری ضروری است \lr{[1{-}5]}. در این فصل، سه محور اصلی حوزه امنیت داده‌های سلامت دیجیتال مورد بررسی و ریشه‌یابی علمی قرار می‌گیرند: (۱) اصول رمزنگاری تصویر مبتنی بر سیستم‌های غیرخطی آشوبی، سیستم‌های ابرآشوبی ۵ بعدی و تبدیل‌های ریاضی، (۲) نقش شبکه‌های عصبی عمیق (به‌ویژه معماری \lr{U{-}Net}) در قطعه‌بندی هوشمند، استخراج ویژگی‌های ناحیه حساس \lr{(ROI)} و تولید کلیدهای رمزی وابسته به محتوا، و (۳) تکنیک‌های پنهان‌نگاری کلاسیک و تطبیقی مبتنی بر برجستگی بصری \lr{(Saliency Detection)} جهت جای‌دهی امن فراداده‌ها بدون ایجاد اعوجاج در کیفیت بالینی \lr{[1, 4, 7]}. بررسی این محورهای سه‌گانه، شالوده نظری و روش‌شناختی لازم برای تحقق معماری هیبریدی پیشنهادی در این پایان‌نامه را فراهم می‌سازد که نمای کلی آن در شکل ۲{-}۱ ترسیم شده است.\par

% SOURCE Introduction_Chapter2.txt:6-6 [spacing]


% SOURCE Introduction_Chapter2.txt:7-7 [image]
\SourcePicture{chapter2_system_architecture.png}{شکل ۲{-}۱: نمودار معماری مفهومی کلان سیستم امنیتی هیبریدی پیشنهادی}{[شکل ۲{-}۱: نمودار معماری مفهومی کلان سیستم امنیتی هیبریدی پیشنهادی {-} \lr{chapter2\_system\_architecture.png}]}{0}

% SOURCE Introduction_Chapter2.txt:8-8 [spacing]


% SOURCE Introduction_Chapter2.txt:9-9 [paragraph]
هدف اصلی از تدوین این فصل، تنها فهرست‌کردن توصیفی یا گزارش ساده نتایج عددی مقالات مرجع نیست؛ بلکه کالبدشکافی تحلیلی هر یک از الگوریتم‌ها، تبیین نحوه فرمول‌بندی ریاضی آن‌ها، ارزیابی شاخص‌های سنجش کارایی (نظیر آنتروپی، \lr{NPCR}، \lr{UACI}، \lr{PSNR} و \lr{SSIM})، و در نهایت استخراج دقیق «شکاف‌های پژوهشی» \lr{(Research Gaps)} موجود در ادبیات تحقیق می‌باشد \lr{[1{-}10]}. بر اساس الگوی استدلال تولمین \lr{(Toulmin Model)}، تحلیل نقادانه کارهای پیشین به عنوان شواهد و مبانی پشتیبان \lr{(Grounds and Backing)} عمل می‌کند تا ضرورت علمی الگوریتم پیشنهادی این پژوهش به عنوان ادعای اصلی \lr{(Claim)} اثبات گردد. کلیه ارقام و شاخص‌های عددی گزارش‌شده در این فصل مستقیماً مستخرج از مقالات مرجع بوده و به عنوان پایه‌های مقایسه‌ای و مبانی سنجش اعتبار پایان‌نامه مورد استفاده قرار گرفته‌اند تا ارزیابی کارایی چارچوب پیشنهادی در فصل چهارم با بالاترین دقت آکادمیک صورت پذیرد.\par

% SOURCE Introduction_Chapter2.txt:10-10 [spacing]


% SOURCE Introduction_Chapter2.txt:11-11 [paragraph]
همچنین این فصل به گونه‌ای سازمان‌دهی شده است که پاسخگوی مستقیم چالش‌های سه‌گانه مطرح‌شده در بیان مسئله فصل اول (کلیدهای رمزی ایستا، همبستگی شدید پیکسلی، و افشای فراداده‌های حساس در بسترهای \lr{IoMT}) باشد. در راستای نیل به این هدف، ساختار فصل دوم در قالب سه بخش اصلی زیر تدوین و هدایت می‌گردد:\par

% SOURCE Introduction_Chapter2.txt:12-12 [spacing]


% SOURCE Introduction_Chapter2.txt:13-13 [paragraph]
۱. \textbf{بخش ۲{-}۲ (مبانی نظری):} تبیین ریاضی و مفهومی اصول رمزنگاری آشوبی، سیستم‌های ۵بعدی، شبکه‌های کانولوشنی عمیق \lr{U{-}Net}، پنهان‌نگاری آگاه به برجستگی دیداری، و فرمول‌های محاسباتی شاخص‌های کمّی ارزیابی امنیت.\par

% SOURCE Introduction_Chapter2.txt:14-14 [paragraph]
۲. \textbf{بخش ۲{-}۳ (پیشینه پژوهش):} مرور نقادانه، دسته‌بندی موضوعی و کالبدشکافی فنی ۱۰ مقاله مرجع برجسته و بین‌المللی (۲۰۲۱ تا ۲۰۲۵) در حوزه رمزنگاری و پنهان‌نگاری تصاویر پزشکی.\par

% SOURCE Introduction_Chapter2.txt:15-15 [paragraph]
۳. \textbf{بخش ۲{-}۴ (ماتریس مقایسه‌ای و استخراج شکاف‌های پژوهشی):} تدوین جدول جامع تاکسونومی پیشینه، استخراج ۳ شکاف پژوهشی اساسی در ادبیات تحقیق، و زمینه‌سازی مستقیم برای ارائه معماری ۴لایه پیشنهادی در فصل سوم.\par

% SOURCE Introduction_Chapter2.txt:16-16 [spacing]


% SOURCE Introduction_Chapter2.txt:17-17 [paragraph]
سیر منطقی، پیوستگی روایی \lr{(Red Thread)} و نقشه راه مفهومی این فصل در شکل ۲{-}۲ به تصویر کشیده شده است.\par

% SOURCE Introduction_Chapter2.txt:18-18 [spacing]


% SOURCE Introduction_Chapter2.txt:19-19 [image]
\SourcePicture{chapter2_roadmap.png}{شکل ۲{-}۲: نمودار نقشه راه مفهومی و ساختاری فصل دوم (از مبانی نظری تا استخراج شکاف‌های پژوهشی)}{[شکل ۲{-}۲: نمودار نقشه راه مفهومی و ساختاری فصل دوم (از مبانی نظری تا استخراج شکاف‌های پژوهشی) {-} \lr{chapter2\_roadmap.png}]}{0}

% SOURCE Introduction_Chapter2.txt:20-20 [spacing]

% SOURCE Theoretical_Foundations_Chapter2.txt:1-1 [heading]
\SourceHeading{1}{۲{-}۲. مبانی نظری}

% SOURCE Theoretical_Foundations_Chapter2.txt:2-2 [spacing]


% SOURCE Theoretical_Foundations_Chapter2.txt:3-3 [heading]
\SourceHeading{2}{۲{-}۲{-}۱. رمزنگاری تصویر دیجیتال و اصول جایگشت و انتشار}

% SOURCE Theoretical_Foundations_Chapter2.txt:4-4 [paragraph]
در رمزنگاری تصویر دیجیتال، داده‌های پیکسل از فضای واضح به نمایش رمزشده تبدیل می‌شوند به‌گونه‌ای که بازیابی اطلاعات اصلی بدون در اختیار داشتن کلید معتبر غیرممکن گردد. تصاویر پزشکی به دلیل دارا بودن ویژگی‌های ذاتی منحصربه‌فرد از جمله حجم سنگین داده، افزونگی اطلاعاتی بسیار بالا و همبستگی شدید پیکسلی میان عناصر همجوار در سه جهت افقی، عمودی و قطری (با ضریب همبستگی $r \approx 1$)، نیازمند معماری‌های رمزی ویژه‌ای هستند؛ زیرا الگوریتم‌های سنتی رمزشده متنی (مانند \lr{AES} و \lr{RSA}) به دلیل سرعت پایین پردازش و عدم توانایی در شکستن همبستگی دو‌بعدی پیکسل‌ها، دچار گلوگاه زمانی می‌شوند \lr{[2, 3]}.\par

% SOURCE Theoretical_Foundations_Chapter2.txt:5-5 [spacing]


% SOURCE Theoretical_Foundations_Chapter2.txt:6-6 [paragraph]
بر اساس الگوی ارائه شده توسط کلود شانون، دو ساختار حیاتی در طراحی سیستم‌های امنیتی رمزنگاری تصویر عبارتند از:\par

% SOURCE Theoretical_Foundations_Chapter2.txt:7-7 [paragraph]
۱. جایگشت یا آشفته‌سازی \lr{(Confusion / Permutation)}: تغییر موقعیت مکانی و جابه‌جایی آدرس پیکسل‌ها در ماتریس تصویر بدون تغییر در مقادیر عددی و شدت روشنایی آن‌ها جهت از بین بردن همبستگی مکانی و ساختاری.\par

% SOURCE Theoretical_Foundations_Chapter2.txt:8-8 [paragraph]
۲. انتشار \lr{(Diffusion)}: تغییر مقادیر عددی و شدت روشنایی پیکسل‌ها به‌گونه‌ای که تغییر در یک پیکسل از تصویر ورودی واضح، منجر به تغییر گسترده و غیرقابل پیش‌بینی در سراسر خروجی رمزشده گردد (محقق ساختن اثر بهمنی کامل یا \lr{Avalanche Effect}).\par

% SOURCE Theoretical_Foundations_Chapter2.txt:9-9 [spacing]


% SOURCE Theoretical_Foundations_Chapter2.txt:10-10 [heading]
\SourceHeading{2}{۲{-}۲{-}۲. سیستم‌های غیرخطی آشوبی و ابرآشوبی}

% SOURCE Theoretical_Foundations_Chapter2.txt:11-11 [paragraph]
سیستم‌های آشوبی، سیستم‌های دینامیکی غیرخطی و معینی هستند که در محدوده مشخصی از پارامترهای کنترلی، رفتاری شبه‌تصادفی، غیردوره‌ای، حساس به شرایط اولیه و غیرقابل پیش‌بینی در درازمدت از خود نشان می‌دهند. عدم قطعیت ذاتی و حساسیت شدید به مقادیر اولیه (معروف به اثر پروانه‌ای یا \lr{Butterfly Effect})، این سیستم‌ها را به ابزاری فوق‌العاده کارآمد برای تولید توالی‌های کلید رمزی و ماتریس‌های آشفته‌سازی تبدیل کرده است \lr{[1, 3]}.\par

% SOURCE Theoretical_Foundations_Chapter2.txt:12-12 [spacing]


% SOURCE Theoretical_Foundations_Chapter2.txt:13-13 [paragraph]
در حالت کلی، سیستم‌های آشوبی کم‌بعد (مانند نگاشت لوجستیک ۱ بعدی یا نگاشت‌های ۲ بعدی ماتیویی و هنون) دارای فضای کلید محدود (کمتر از $2^{128}$) بوده و به دلیل سادگی ساختار ریاضی، در برابر حملات بازسازی فضای فاز \lr{(Phase Space Reconstruction)} و حملات انتخاب متن‌واضح \lr{(CPA)} آسیب‌پذیرند. برای غلبه بر این محدودیت بنیادی، سیستم‌های ابرآشوبی \lr{(Hyperchaotic Systems)} با حداقل ۴ بعد در فضای فاز و دارا بودن بیش از یک توان نمای لیاپانوف مثبت \lr{(Multiple Positive Lyapunov Exponents)} توسعه یافته‌اند. وجود چندین توان لیاپانوف مثبت باعث می‌شود که توالی‌های تولیدشده دارای پیچیدگی دینامیکی فوق‌العاده بالا، آنتروپی آرمانی و فضای کلید فراتر از $2^{500}$ باشند که هرگونه حمله جستجوی فراگیر \lr{(Brute{-}force)} را تماماً منتفی می‌سازد \lr{[1]}.\par

% SOURCE Theoretical_Foundations_Chapter2.txt:14-14 [spacing]


% SOURCE Theoretical_Foundations_Chapter2.txt:15-15 [heading]
\SourceHeading{3}{۲{-}۲{-}۲{-}۱. فرمول‌بندی ریاضی نگاشت‌های آشوبی و سیستم ۵ بعدی}

% SOURCE Theoretical_Foundations_Chapter2.txt:16-16 [paragraph]
الف) نگاشت لوجستیک ۱ بعدی \lr{(1D Logistic Map)}:\par

% SOURCE Theoretical_Foundations_Chapter2.txt:17-17 [paragraph]
این نگاشت غیرخطی با معادله دیفرانسیل گسسته زیر تعریف می‌شود \lr{[3]}:\par

% SOURCE Theoretical_Foundations_Chapter2.txt:18-18 [spacing]


% SOURCE Theoretical_Foundations_Chapter2.txt:19-19 [image]
\SourcePicture{eq2_1_logistic_map.png}{رابطه ۲{-}۱: معادله دیفرانسیل گسسته نگاشت لوجستیک ۱ بعدی}{[رابطه ۲{-}۱: معادله دیفرانسیل گسسته نگاشت لوجستیک ۱ بعدی {-} \lr{eq2\_1\_logistic\_map.png}]}{2}

% SOURCE Theoretical_Foundations_Chapter2.txt:20-20 [spacing]


% SOURCE Theoretical_Foundations_Chapter2.txt:21-21 [paragraph]
که در آن $x_n \in (0,\allowbreak{} 1)$ متغیر حالت سیستم در گام $n$ و $\mu \in (0,\allowbreak{} 4.0]$ پارامتر کنترل یا انشعاب \lr{(Bifurcation Parameter)} است. سیستم در محدوده $3.57 < \mu \le 4.0$ رفتار کاملاً آشوبی از خود نشان می‌دهد.\par

% SOURCE Theoretical_Foundations_Chapter2.txt:22-22 [spacing]


% SOURCE Theoretical_Foundations_Chapter2.txt:23-23 [paragraph]
ب) سیستم ابرآشوبی ۵ بعدی پیشنهادی سوباترا \lr{(5D Hyperchaotic System)}:\par

% SOURCE Theoretical_Foundations_Chapter2.txt:24-24 [paragraph]
معادلات دیفرانسیل حالت سیستم ابرآشوبی ۵ بعدی پیوسته به صورت فرمول‌بندی غیرخطی زیر بیان می‌گردند \lr{[1]}:\par

% SOURCE Theoretical_Foundations_Chapter2.txt:25-25 [spacing]


% SOURCE Theoretical_Foundations_Chapter2.txt:26-26 [image]
\SourcePicture{eq2_2_5d_hyperchaos.png}{رابطه ۲{-}۲: دستگاه معادلات دیفرانسیل غیرخطی سیستم ابرآشوبی ۵ بعدی سوباترا}{[رابطه ۲{-}۲: دستگاه معادلات دیفرانسیل غیرخطی سیستم ابرآشوبی ۵ بعدی سوباترا {-} \lr{eq2\_2\_5d\_hyperchaos.png}]}{2}

% SOURCE Theoretical_Foundations_Chapter2.txt:27-27 [spacing]


% SOURCE Theoretical_Foundations_Chapter2.txt:28-28 [paragraph]
که در آن $(x,\allowbreak{} y,\allowbreak{} z,\allowbreak{} u,\allowbreak{} v)$ متغیرهای حالت سیستم در فضای پنج‌بعدی و $(a,\allowbreak{} b,\allowbreak{} c,\allowbreak{} d,\allowbreak{} k,\allowbreak{} h,\allowbreak{} w)$ پارامترهای کنترلی غیرخطی سیستم می‌باشند. محاسبه توان‌های نمای لیاپانوف $\lambda_1,\allowbreak{} \lambda_2 > 0$ وجود حداقل دو توان مثبت را اثبات کرده و رفتار ابرآشوبی غریب \lr{(Strange Attractor)} سیستم را تایید می‌نماید.\par

% SOURCE Theoretical_Foundations_Chapter2.txt:29-29 [spacing]


% SOURCE Theoretical_Foundations_Chapter2.txt:30-30 [heading]
\SourceHeading{2}{۲{-}۲{-}۳. یادگیری عمیق در امنیت تصاویر پزشکی}

% SOURCE Theoretical_Foundations_Chapter2.txt:31-31 [paragraph]
در سال‌های اخیر، استفاده از شبکه‌های عصبی عمیق \lr{(Deep Neural Networks)} از نقش سنتی آن‌ها در تشخیص و دسته‌بندی بیماری‌ها فراتر رفته و به حوزه‌های رمزنگاری و پنهان‌نگاری گسترش یافته است. شبکه عصبی می‌تواند به عنوان یک ماژول تصمیم‌گیری تطبیقی، استخراج‌کننده خصوصیات آماری، یا الگوریتم قطعه‌بندی معنایی نواحی حساس \lr{(ROI)} عمل نماید \lr{[1, 2, 7]}.\par

% SOURCE Theoretical_Foundations_Chapter2.txt:32-32 [spacing]


% SOURCE Theoretical_Foundations_Chapter2.txt:33-33 [heading]
\SourceHeading{3}{۲{-}۲{-}۳{-}۱. فرمول‌بندی ریاضی و معماری شبکه \lr{U{-}Net}}

% SOURCE Theoretical_Foundations_Chapter2.txt:34-34 [paragraph]
معماری \lr{U{-}Net} یک شبکه کانولوشنی متقارن شامل دو مسیر اصلی است:\par

% SOURCE Theoretical_Foundations_Chapter2.txt:35-35 [paragraph]
۱. مسیر انقباضی \lr{(Encoder / Contracting Path)}: برای استخراج ویژگی‌های معنایی در سطوح مختلف از طریق اعمال لایه‌های کانولوشن دو‌بعدی، تابع فعال‌سازی \lr{ReLU} و لایه‌های کاهش بعد \lr{Max{-}Pooling} (با اندازه $2 \times 2$ و گام ۲).\par

% SOURCE Theoretical_Foundations_Chapter2.txt:36-36 [paragraph]
۲. مسیر انبساطی \lr{(Decoder / Expanding Path)}: برای بازیابی ابعاد مکانی و بازسازی ماسک دقت بالا از طریق کانولوشن معکوس \lr{(Up{-}convolution)} و استفاده از اتصالات پرشی \lr{(Skip Connections)}. اتصالات پرشی ویژگی‌های با وضوح بالای لایه‌های \lr{Encoder} را مستقیماً به لایه‌های \lr{Decoder} متصل می‌کنند.\par

% SOURCE Theoretical_Foundations_Chapter2.txt:37-37 [spacing]


% SOURCE Theoretical_Foundations_Chapter2.txt:38-38 [paragraph]
فرمول‌بندی لایه‌های کانولوشن، ادغام اتصالات پرشی و تابع زیان \lr{(Loss Function)} ترکیبی شبکه \lr{U{-}Net} به صورت زیر تعریف می‌گردند \lr{[1]}:\par

% SOURCE Theoretical_Foundations_Chapter2.txt:39-39 [spacing]


% SOURCE Theoretical_Foundations_Chapter2.txt:40-40 [image]
\SourcePicture{eq2_3_unet_math.png}{رابطه ۲{-}۳: فرمول‌بندی لایه‌های کانولوشن، ادغام \lr{Skip Connection} و تابع زیان ترکیبی \lr{U{-}Net}}{[رابطه ۲{-}۳: فرمول‌بندی لایه‌های کانولوشن، ادغام \lr{Skip Connection} و تابع زیان ترکیبی \lr{U{-}Net {-} eq2\_3\_unet\_math.png}]}{2}

% SOURCE Theoretical_Foundations_Chapter2.txt:41-41 [spacing]


% SOURCE Theoretical_Foundations_Chapter2.txt:42-42 [paragraph]
تابع زیان ترکیبی $\mathcal{L}_{\text{total}}$ با تلفیق انتروپی متقابل دودویی \lr{(BCE)} و ضریب دایس \lr{(Dice Loss)}، قطعیت قطعه‌بندی دقیق ناحیه \lr{ROI} را بر روی مدالیته‌های مختلف پزشکی تضمین می‌نماید.\par

% SOURCE Theoretical_Foundations_Chapter2.txt:43-43 [spacing]


% SOURCE Theoretical_Foundations_Chapter2.txt:44-44 [heading]
\SourceHeading{2}{۲{-}۲{-}۴. پنهان‌نگاری تطبیقی و تحلیل برجستگی بصری \lr{(Saliency Detection)}}

% SOURCE Theoretical_Foundations_Chapter2.txt:45-45 [paragraph]
پنهان‌نگاری \lr{(Steganography)} علم و هنر پنهان‌سازی اطلاعات محرمانه (فراداده‌های هویتی و بالینی بیمار) در درون یک تصویر پوششی پزشکی است به‌گونه‌ای که تغییرات ایجادشده برای سیستم بصری انسان \lr{(HVS)} و ابزارهای پنهان‌کاوی \lr{(Steganalysis)} غیرقابل تشخیص باشد. الگوریتم‌های سنتی پنهان‌نگاری (نظیر \lr{LSB} و \lr{PVD}) داده‌ها را به صورت غیرتطبیقی در پیکسل‌ها جای‌دهی می‌کنند که این امر منجر به ایجاد نویز دیداری و افت کیفیت تشخیصی در ناحیه حساس \lr{ROI} می‌گردد \lr{[5]}.\par

% SOURCE Theoretical_Foundations_Chapter2.txt:46-46 [spacing]


% SOURCE Theoretical_Foundations_Chapter2.txt:47-47 [paragraph]
تکنیک‌های نوین پنهان‌نگاری تطبیقی بر پایه نقشه برجستگی بصری \lr{(Saliency Map)} عمل می‌کنند. نقشه برجستگی با تحلیل درجه توجه دیداری، پیکسل‌های تصویر را به دو دسته نواحی برجسته \lr{(High Saliency / ROI)} و نواحی کم‌برجستگی پس‌زمینه \lr{(Low Saliency / Background)} تقسیم می‌کند. در این روش، فراداده متنی بیمار ابتدا توسط الگوریتم‌های فشرده‌سازی بدون اتلاف (مانند \lr{zlib}) فشرده شده و سپس صرفاً در پیکسل‌های پس‌زمینه جای‌دهی می‌شود تا \lr{PSNR} بالاتر از $50\text{ dB}$ حفظ شده و ناحیه تشخیصی \lr{ROI} ۱۰۰\SourceGlyph{066A} دست‌نخورده باقی بماند \lr{[4, 5]}.\par

% SOURCE Theoretical_Foundations_Chapter2.txt:48-48 [spacing]


% SOURCE Theoretical_Foundations_Chapter2.txt:49-49 [image]
\SourcePicture{chapter2_system_architecture.png}{شکل ۲{-}۳: نمودار معماری مفهومی زنجیره پردازشی امنیت هیبریدی تصویر پزشکی}{[شکل ۲{-}۳: نمودار معماری مفهومی زنجیره پردازشی امنیت هیبریدی تصویر پزشکی {-} \lr{chapter2\_system\_architecture.png}]}{0}

% SOURCE Theoretical_Foundations_Chapter2.txt:50-50 [spacing]


% SOURCE Theoretical_Foundations_Chapter2.txt:51-51 [heading]
\SourceHeading{2}{۲{-}۲{-}۵. فرمول‌های ریاضی و شاخص‌های ارزیابی کمّی امنیت و کیفیت}

% SOURCE Theoretical_Foundations_Chapter2.txt:52-52 [paragraph]
برای ارزیابی کمّی کارایی الگوریتم‌های رمزنگاری و پنهان‌نگاری، شاخص‌های ریاضی زیر در ادبیات پژوهش تعریف و استانداردسازی شده‌اند:\par

% SOURCE Theoretical_Foundations_Chapter2.txt:53-53 [spacing]


% SOURCE Theoretical_Foundations_Chapter2.txt:54-54 [paragraph]
الف) آنتروپی شانون \lr{(Shannon Entropy)} و ضریب همبستگی پیکسل‌های مجاور \lr{(Correlation Coefficient)}:\par

% SOURCE Theoretical_Foundations_Chapter2.txt:55-55 [paragraph]
آنتروپی معیاری برای سنجش درجه تصادفی بودن خروجی رمزشده است و ضریب همبستگی میزان وابستگی مکانی پیکسل‌های همجوار را می‌سنجد \lr{[1, 3]}:\par

% SOURCE Theoretical_Foundations_Chapter2.txt:56-56 [spacing]


% SOURCE Theoretical_Foundations_Chapter2.txt:57-57 [image]
\SourcePicture{eq2_4_entropy_correlation.png}{رابطه ۲{-}۴: فرمول‌های آنتروپی اطلاعات شانون و ضریب همبستگی پیکسل‌های مجاور}{[رابطه ۲{-}۴: فرمول‌های آنتروپی اطلاعات شانون و ضریب همبستگی پیکسل‌های مجاور {-} \lr{eq2\_4\_entropy\_correlation.png}]}{2}

% SOURCE Theoretical_Foundations_Chapter2.txt:58-58 [spacing]


% SOURCE Theoretical_Foundations_Chapter2.txt:59-59 [paragraph]
برای یک تصویر رمزشده ۸ بیتی خاکستری، مقدار آرمانی آنتروپی برابر $8.0$ بیت و ضریب همبستگی آرمانی $r_{xy} \approx 0.0$ است.\par

% SOURCE Theoretical_Foundations_Chapter2.txt:60-60 [spacing]


% SOURCE Theoretical_Foundations_Chapter2.txt:61-61 [paragraph]
ب) شاخص‌های ارزیابی حملات تفاضلی (\lr{NPCR} و \lr{UACI}):\par

% SOURCE Theoretical_Foundations_Chapter2.txt:62-62 [paragraph]
این دو معیار، میزان حساسیت تصویر رمزشده به تغییر یک پیکسلی در تصویر واضح ورودی را ارزیابی می‌کنند \lr{[1, 3]}:\par

% SOURCE Theoretical_Foundations_Chapter2.txt:63-63 [spacing]


% SOURCE Theoretical_Foundations_Chapter2.txt:64-64 [image]
\SourcePicture{eq2_5_npcr_uaci.png}{رابطه ۲{-}۵: فرمول‌های معیارهای حساسیت حملات تفاضلی (\lr{NPCR} و \lr{UACI})}{[رابطه ۲{-}۵: فرمول‌های معیارهای حساسیت حملات تفاضلی (\lr{NPCR} و \lr{UACI}) {-} \lr{eq2\_5\_npcr\_uaci.png}]}{2}

% SOURCE Theoretical_Foundations_Chapter2.txt:65-65 [spacing]


% SOURCE Theoretical_Foundations_Chapter2.txt:66-66 [paragraph]
مقادیر آرمانی استاندارد برای \lr{NPCR} بالاتر از $99.60\%$ و برای \lr{UACI} حدود $33.46\%$ می‌باشد.\par

% SOURCE Theoretical_Foundations_Chapter2.txt:67-67 [spacing]


% SOURCE Theoretical_Foundations_Chapter2.txt:68-68 [paragraph]
ج) معیارهای کیفیت دیداری پنهان‌نگاری (\lr{PSNR} و \lr{SSIM}):\par

% SOURCE Theoretical_Foundations_Chapter2.txt:69-69 [paragraph]
برای ارزیابی عدم اعوجاج و میزان شفافیت تصویر پنهان‌نگاری‌شده نسبت به تصویر اصلی، از نسبت سیگنال به نویز بیشینه \lr{(PSNR)} و شاخص شباهت ساختاری \lr{(SSIM)} استفاده می‌شود \lr{[4, 7]}:\par

% SOURCE Theoretical_Foundations_Chapter2.txt:70-70 [spacing]


% SOURCE Theoretical_Foundations_Chapter2.txt:71-71 [image]
\SourcePicture{eq2_6_psnr_ssim.png}{رابطه ۲{-}۶: فرمول‌های شاخص کیفیت دیداری \lr{(PSNR)} و شباهت ساختاری \lr{(SSIM)}}{[رابطه ۲{-}۶: فرمول‌های شاخص کیفیت دیداری \lr{(PSNR)} و شباهت ساختاری \lr{(SSIM)} {-} \lr{eq2\_6\_psnr\_ssim.png}]}{2}

% SOURCE Theoretical_Foundations_Chapter2.txt:72-72 [spacing]


% SOURCE Theoretical_Foundations_Chapter2.txt:73-73 [paragraph]
مقدار \lr{PSNR} بالای $50\text{ dB}$ و \lr{SSIM} نزدیک به $1.0$ نشان‌دهنده حفظ کامل کیفیت دیداری و عدم اعوجاج در تصویر پزشکی است.\par

% SOURCE Theoretical_Foundations_Chapter2.txt:74-74 [spacing]

% SOURCE Literature_Review_Chapter2.txt:1-1 [heading]
\SourceHeading{1}{۲{-}۳. بررسی پیشینه پژوهش (مرور مقالات مرجع)}

% SOURCE Literature_Review_Chapter2.txt:2-2 [spacing]


% SOURCE Literature_Review_Chapter2.txt:3-3 [paragraph]
برای ارزیابی جامع و نقادانه کارهای پیشین، تبیین نحوه تکامل الگوریتم‌های امنیتی و استخراج دقیق شکاف‌های پژوهشی \lr{(Research Gaps)}، ۱۰ مقاله مرجع برجسته و محوری موجود در ادبیات تحقیق (منتشرشده بین سال‌های ۲۰۲۱ تا ۲۰۲۵) بر پایه ۳ خوشه اصلی تکنولوژیک دسته‌بندی و تحلیل می‌گردند. سیر دسته‌بندی موضوعی مقالات و ارتباط آن‌ها با ارکان نظام امنیتی پیشنهادی در شکل ۲{-}۴ نشان داده شده است.\par

% SOURCE Literature_Review_Chapter2.txt:4-4 [spacing]


% SOURCE Literature_Review_Chapter2.txt:5-5 [image]
\SourcePicture{chapter2_literature_classification.png}{شکل ۲{-}۴: نمودار درخت دسته‌بندی پیشینه پژوهش بر اساس سه محور اصلی تکنولوژیک}{[شکل ۲{-}۴: نمودار درخت دسته‌بندی پیشینه پژوهش بر اساس سه محور اصلی تکنولوژیک {-} \lr{chapter2\_literature\_classification.png}]}{0}

% SOURCE Literature_Review_Chapter2.txt:6-6 [spacing]


% SOURCE Literature_Review_Chapter2.txt:7-7 [heading]
\SourceHeading{2}{۲{-}۳{-}۱. سیستم‌های غیرخطی آشوبی و ابرآشوبی در رمزنگاری تصویر پزشکی}

% SOURCE Literature_Review_Chapter2.txt:8-8 [spacing]


% SOURCE Literature_Review_Chapter2.txt:9-9 [paragraph]
در این دسته از پژوهش‌ها، تمرکز بر به‌کارگیری خواص غیرخطی نگاشت‌های آشوبی جهت ایجاد آشفته‌سازی \lr{(Confusion)} و انتشار \lr{(Diffusion)} در ماتریس پیکسل‌های تصویر استوار است:\par

% SOURCE Literature_Review_Chapter2.txt:10-10 [spacing]


% SOURCE Literature_Review_Chapter2.txt:11-11 [paragraph]
۱. بررسی مقاله \lr{Pourasad et al}. \lr{(2021)} {-} ترکیب نگاشت‌های آشوبی و حوزه تبدیل موجک \lr{[3]}:\par

% SOURCE Literature_Review_Chapter2.txt:12-12 [paragraph]
پوراسد و همکاران الگوریتمی دو مرحله‌ای برای رمزنگاری تصاویر دیجیتال بر پایه ترکیب نگاشت‌های آشوبی لوجستیک ۱ بعدی و تبدیل موجک گسسته \lr{(DWT)} ارائه دادند. در این طرح، ابتدا عملیات انتشار توسط نگاشت لوجستیک بر روی تصویر واضح اعمال می‌شود؛ سپس با تجزیه موجکی، زیرباندهای فرکانسی \lr{(LL, LH, HL, HH)} استخراج گردیده و عملیات جایگشت آدرس پیکسل‌ها در حوزه فرکانس توسط شبکه نقشه‌های جفت‌شده \lr{(CML)} انجام می‌گیرد. نتایج عددی گزارش‌شده برای تصاویر معیار شامل آنتروپی شانون \lr{7.9978}، شاخص \lr{NPCR} در محدوده ۹۹.۶۲۵\SourceGlyph{066A} تا ۹۹.۸۸۱\SourceGlyph{066A} و \lr{UACI} در محدوده ۳۳.۱۲۰\SourceGlyph{066A} تا ۳۳.۶۷۱\SourceGlyph{066A} می‌باشد. با این حال، نقد اصلی داوری بر این مقاله، آزمون الگوریتم بر روی تصاویر عمومی (مانند \lr{Lena}) و عدم بررسی دادگان حساس پزشکی، و همچنین حساسیت شدید حوزه تبدیل \lr{DWT} به خطاهای گردکردن اعداد اعشاری هنگام بازسازی دقیق تصویر تشخیصی است.\par

% SOURCE Literature_Review_Chapter2.txt:13-13 [spacing]


% SOURCE Literature_Review_Chapter2.txt:14-14 [paragraph]
۲. بررسی مقاله \lr{Subathra \& Thanikaiselvan} \lr{(2025)} {-} قطعه‌بندی \lr{U{-}Net} و ابرآشوب ۵ بعدی \lr{[1]}:\par

% SOURCE Literature_Review_Chapter2.txt:15-15 [paragraph]
سوباترا و تانیکایسلوان چارچوبی امنیتی و هوشمند شامل شبکه عصبی عمیق \lr{U{-}Net} جهت قطعه‌بندی ناحیه حساس تشخیصی \lr{(ROI)}، سیستم ابرآشوبی ۵ بعدی پیوسته و فازهای انتشار پویای \lr{DNA} ارائه کردند. ویژگی‌های آماری \lr{ROI} (شامل مجموع شدت روشنایی، کجی و کشیدگی هیستوگرام) برای مقداردهی پویای شرایط اولیه سیستم ۵ بعدی استفاده می‌شوند. الگوریتم شامل جایگشت زیگ‌زاگی و انتشار پویای \lr{DNA} \lr{(Dynamic DNA Flip / XOR)} است. نتایج عددی گزارش‌شده نشان‌دهنده آنتروپی شانون \lr{7.9971}، \lr{NPCR} بالای \lr{99.61\%}، \lr{UACI} حدود \lr{33.49\%} و زمان اجرای ۲.۹۳ ثانیه برای تصاویر \lr{\$256 \SourceControl{0009}imes 256\$} است. نقد وارد بر این طرح، عدم وجود ماژول پنهان‌نگاری فراداده بیمار، عدم حفاظت از اطلاعات هویتی و عدم ارائه مکانیزم امن برای انتقال شرایط اولیه تولیدشده به سمت گیرنده است.\par

% SOURCE Literature_Review_Chapter2.txt:16-16 [spacing]


% SOURCE Literature_Review_Chapter2.txt:17-17 [paragraph]
۳. بررسی مقاله \lr{Preethi et al}. \lr{(2025)} {-} سیستم‌های رمزی ناهمگون آشوبی \lr{[10]}:\par

% SOURCE Literature_Review_Chapter2.txt:18-18 [paragraph]
پریتی و همکاران یک سیستم رمزنگاری ناهمگون \lr{(Asymmetric Chaotic Cryptosystem)} ترکیبی از زیرساخت کلید عمومی{-}خصوصی \lr{(PKI)} و سیستم‌های آشوبی برای تصاویر پزشکی \lr{MRI} و \lr{CT} توسعه دادند. این مقاله با استفاده از کلیدهای ناهمگون، مشکل تبادل امن کلید در شبکه‌های عمومی را حل می‌نماید. نتایج شبیه‌سازی در محیط متلب مقاومت مناسب در برابر حملات تحلیل کلید و آنتروپی بالای \lr{7.996} بیت را تایید می‌کند؛ اما نقد اصلی بر این پژوهش، عدم به‌کارگیری قطعه‌بندی هوشمند \lr{ROI}، عدم تولید کلید وابسته به تصویر، و عدم پشتیبانی از پنهان‌نگاری فراداده‌های هویتی بیمار می‌باشد.\par

% SOURCE Literature_Review_Chapter2.txt:19-19 [spacing]


% SOURCE Literature_Review_Chapter2.txt:20-20 [heading]
\SourceHeading{2}{۲{-}۳{-}۲. شبکه‌های عصبی عمیق در استخراج ویژگی، قطعه‌بندی و رمزنگاری}

% SOURCE Literature_Review_Chapter2.txt:21-21 [spacing]


% SOURCE Literature_Review_Chapter2.txt:22-22 [paragraph]
در این گروه از مقالات، از هوش مصنوعی و یادگیری عمیق جهت هوشمندسازی فرایند رمزنگاری، تولید کلید پویا و جداسازی نواحی حساس استفاده شده است:\par

% SOURCE Literature_Review_Chapter2.txt:23-23 [spacing]


% SOURCE Literature_Review_Chapter2.txt:24-24 [paragraph]
۴. بررسی مقاله \lr{Alarood et al}. \lr{(2023)} {-} عمیق‌سازی رمزنگاری \lr{E{-}Health} با شبکه \lr{DNN} و نگاشت هنون \lr{[2]}:\par

% SOURCE Literature_Review_Chapter2.txt:25-25 [paragraph]
الارود و همکاران سیستمی برای انتقال امن تصاویر پزشکی در بسترهای \lr{E{-}Health} با استفاده از شبکه‌های عصبی عمیق \lr{(DNN)} و نگاشت آشوبی ۲ بعدی هنون ارائه دادند. در این روش، وزن‌های آموزش‌دیده شبکه عصبی جهت تقسیم هوشمند تصویر به بلوک‌های \lr{ROI} و غیر \lr{ROI} و تولید توالی‌های آشفته‌سازی به کار می‌روند. نتایج گزارش‌شده شامل \lr{NPCR} حدود ۹۹.۶۱\SourceGlyph{066A}، \lr{UACI} حدود ۳۳.۴۲\SourceGlyph{066A} و ضریب همبستگی بین {-}۰.۰۴۱ تا ۰.۰۸۱ است. محدودیت اصلی این طرح، بار محاسباتی بسیار سنگین آموزش شبکه عمیق، استفاده از کلیدهای رمزی ثابت و عدم مقاومت در برابر حملات فشرده‌سازی و تغییر اندازه تصویر است.\par

% SOURCE Literature_Review_Chapter2.txt:26-26 [spacing]


% SOURCE Literature_Review_Chapter2.txt:27-27 [paragraph]
۵. بررسی مقاله \lr{Husamuldeen et al}. \lr{(2025)} {-} تولید کلید پویا با \lr{EXIF Metadata} و شبکه‌های \lr{ResNet{-}18 / MobileNetV2} \lr{[8]}:\par

% SOURCE Literature_Review_Chapter2.txt:28-28 [paragraph]
حسام‌الدین و همکاران سیستمی برای رمزنگاری تصاویر پزشکی ارائه دادند که فراداده‌های \lr{EXIF} تصویر را استخراج کرده و از طریق شبکه‌های عمیق \lr{ResNet{-}18} و \lr{MobileNetV2} برای تولید کلیدهای پویا و تنظیم پارامترهای کنترل سیستم‌های آشوبی به کار می‌گیرد. نتایج شامل آنتروپی \lr{7.9988} و \lr{NPCR} بالای \lr{99.62\%} است. نقد وارد بر این پژوهش، عدم حفاظت از خود فراداده‌های \lr{EXIF} هنگام انتقال بر روی شبکه و مخاطره جدّی افشای اطلاعات هویتی بیمار در صورت استراق سمع نشت داده است.\par

% SOURCE Literature_Review_Chapter2.txt:29-29 [spacing]


% SOURCE Literature_Review_Chapter2.txt:30-30 [paragraph]
۶. بررسی مقاله \lr{Kumar et al}. \lr{(2025)} {-} رمزنگاری عمیق با \lr{DCGAN} و دامنه سیاره مجازی \lr{(VPD)} \lr{[9]}:\par

% SOURCE Literature_Review_Chapter2.txt:31-31 [paragraph]
کومار و همکاران طرحی هیبریدی مبتنی بر شبکه‌های مولد تصادفی عمیق \lr{(DCGAN)}، دامنه سیاره مجازی \lr{(Virtual Planet Domain)} و نگاشت‌های غیرخطی \lr{1{-}DEC} ارائه کردند. در این طرح، مولد \lr{(Generator)} شبکه عمیق، تصاویر فریبنده \lr{(Decoy Images)} تولید می‌کند تا الگوی رمزی اصلی از چشم مهاجمان پنهان بماند. آزمون‌های استاندارد \lr{NIST} مقاومت بالای طرح را تایید می‌کنند؛ اما نقد داوری بر این مقاله، زمان اجرای بسیار سنگین، عدم امکان پیاده‌سازی زمان‌واقعی \lr{(Real{-}time)} و بار پردازشی بالایی است که به تجهیزات کم‌مصرف \lr{IoMT} تحمیل می‌نماید.\par

% SOURCE Literature_Review_Chapter2.txt:32-32 [spacing]


% SOURCE Literature_Review_Chapter2.txt:33-33 [heading]
\SourceHeading{2}{۲{-}۳{-}۳. تکنیک‌های پنهان‌نگاری تطبیقی و حفظ حریم خصوصی بیماران}

% SOURCE Literature_Review_Chapter2.txt:34-34 [spacing]


% SOURCE Literature_Review_Chapter2.txt:35-35 [paragraph]
این دسته از پژوهش‌ها بر پنهان‌سازی فراداده‌ها و پرونده‌های بالینی بیمار در درون تصویر پزشکی بدون ایجاد اعوجاج تشخیصی تمرکز دارند:\par

% SOURCE Literature_Review_Chapter2.txt:36-36 [spacing]


% SOURCE Literature_Review_Chapter2.txt:37-37 [paragraph]
۷. بررسی مقاله \lr{Y\SourceGlyph{0131}ld\SourceGlyph{0131}r\SourceGlyph{0131}m} \lr{(2021)} {-} پنهان‌سازی صوت رمزنگاری‌شده در تصاویر \lr{CT} کووید{-}۱۹ \lr{[4]}:\par

% SOURCE Literature_Review_Chapter2.txt:38-38 [paragraph]
ییلدیریم روشی برای پنهان‌سازی نظرات صوتی رمزنگاری‌شده پزشکان در تصاویر \lr{CT} ریه بیماران مبتلا به کرونا ارائه داد. صوت پزشک ابتدا به فرمت تصویری تبدیل شده، با نگاشت‌های آشوبی رمزنگاری گردیده و سپس در بیت‌های کم‌ارزش \lr{(LSB)} تصویر پزشکی جای‌دهی می‌شود. نتایج عددی گزارش‌شده شامل آنتروپی \lr{7.9993}، \lr{NPCR} برابر \lr{99.8069\%}، \lr{UACI} برابر \lr{33.5688\%} و \lr{SSIM} در کانال‌های رنگی بین \lr{0.7926} تا \lr{0.9273} می‌باشد. نقد داور بر این مقاله، افت معنادار شاخص شباهت ساختاری \lr{(SSIM)} در برخی کانال‌ها به دلیل بار بالای جای‌دهی و عدم تعمیم‌پذیری الگوریتم به سایر مدالیته‌های تشخیصی نظیر \lr{MRI} یا شبکیه چشم است.\par

% SOURCE Literature_Review_Chapter2.txt:39-39 [spacing]


% SOURCE Literature_Review_Chapter2.txt:40-40 [paragraph]
۸. بررسی مقاله \lr{Alshamrani et al}. \lr{(2025)} {-} سامانه \lr{VisionStego AI} و ارزیابی تطبیقی \lr{[5]}:\par

% SOURCE Literature_Review_Chapter2.txt:41-41 [paragraph]
الشمرانی و همکاران سری آزمایش‌های مقایسه‌ای بر روی روش‌های \lr{LSB} و \lr{PVD} انجام داده و سیستم \lr{VisionStego AI} را بر پایه تشخیص برجستگی بصری \lr{(Saliency Detection)} معرفی کردند. در این سیستم، داده‌های فشرده‌سازی شده بیمار با الگوریتم \lr{zlib} صرفاً در نواحی کم‌برجستگی پس‌زمینه جای‌دهی می‌شوند. نتایج نشان‌دهنده تعادل مناسب میان \lr{PSNR} و زمان پردازش است. نقد وارد بر این مقاله، عدم تلفیق ماژول پنهان‌نگاری با الگوریتم‌های رمزنگاری پیشرفته آشوبی و وابستگی شدید کیفیت تصویر به حجم داده‌های جای‌دهی‌شده است.\par

% SOURCE Literature_Review_Chapter2.txt:42-42 [spacing]


% SOURCE Literature_Review_Chapter2.txt:43-43 [paragraph]
۹. بررسی مقاله \lr{Issac et al}. \lr{(2025)} {-} پنهان‌نگاری عمیق با ماژول‌های \lr{SE} و \lr{Inception} بر روی سخت‌افزار لبه \lr{[6]}:\par

% SOURCE Literature_Review_Chapter2.txt:44-44 [paragraph]
ایساک و همکاران معماری پنهان‌نگاری عمیق بر پایه ماژول‌های \lr{Squeeze{-}and{-}Excitation} \lr{(SE)}، \lr{Inception} و اتصالات باقیمانده را برای انتقال داده‌های بزرگ پزشکی بر روی سخت‌افزار لبه \lr{NVIDIA Jetson TX2} ارائه کردند. نتایج گزارش‌شده برای دادگان \lr{MRI} و \lr{OCT} شامل \lr{PSNR} برابر \lr{39.02 dB} و \lr{SSIM} برابر \lr{0.9757} است. نقد اصلی داوری، محدود بودن \lr{PSNR} به زیر ۴۰ دسی‌بل (که در تصاویر پزشکی حساس افت کیفیت بالینی محسوب می‌شود) و پیچیدگی بسیار بالای فاز آموزش شبکه است.\par

% SOURCE Literature_Review_Chapter2.txt:45-45 [spacing]


% SOURCE Literature_Review_Chapter2.txt:46-46 [paragraph]
۱۰. بررسی مقاله \lr{Saidi et al}. \lr{(2024)} {-} پنهان‌نگاری با ظرفیت بالا مبتنی بر \lr{Mask{-}RCNN} و تبدیل \lr{DCT} \lr{[7]}:\par

% SOURCE Literature_Review_Chapter2.txt:47-47 [paragraph]
سعیدی و همکاران روشی برای پنهان‌نگاری با ظرفیت بالا در تصاویر پزشکی ارائه دادند که از شبکه \lr{Mask{-}RCNN} جهت شناسایی هوشمند نواحی پس‌زمینه بی‌اهمیت و تبدیل کوسینوسی گسسته \lr{(DCT)} برای جای‌دهی کد \lr{QR} فراداده بیمار استفاده می‌کند. نتایج آزمایش بر روی دادگان \lr{CHAOS} نشان‌دهنده کیفیت دیداری مناسب است؛ اما نقد وارد بر آن، پیچیدگی زمانی بالای الگوریتم \lr{Mask{-}RCNN} و عدم رمزنگاری آشوبی تصویر اصلی است.\par

% SOURCE Literature_Review_Chapter2.txt:48-48 [spacing]


% SOURCE Literature_Review_Chapter2.txt:49-49 [paragraph]
جدول ۲{-}۱ خلاصه‌ای از ارزیابی کمّی، مدالیته داده‌ها، روش‌های اصلی و محدودیت‌های ۱۰ مقاله مرجع بررسی‌شده را به صورت ماتریس مقایسه‌ای ارائه می‌دهد.\par

% SOURCE Literature_Review_Chapter2.txt:50-50 [spacing]


% SOURCE Literature_Review_Chapter2.txt:51-51 [image]
\SourcePicture{chapter2_taxonomy_table.png}{جدول ۲{-}۱: ماتریس جامع مقایسه‌ای کمّی و ارزیابی نقادانه پیشینه پژوهش}{[جدول ۲{-}۱: ماتریس جامع مقایسه‌ای کمّی و ارزیابی نقادانه پیشینه پژوهش {-} \lr{chapter2\_taxonomy\_table.png}]}{1}

% SOURCE Literature_Review_Chapter2.txt:52-52 [spacing]

% SOURCE Synthesis_Research_Gap_Chapter2.txt:1-1 [heading]
\SourceHeading{1}{۲{-}۴. جمع‌بندی تحلیلی، ماتریس مقایسه‌ای و استخراج شکاف پژوهشی}

% SOURCE Synthesis_Research_Gap_Chapter2.txt:2-2 [spacing]


% SOURCE Synthesis_Research_Gap_Chapter2.txt:3-3 [heading]
\SourceHeading{2}{۲{-}۴{-}۱. ماتریس جامع مقایسه‌ای کارهای پیشین (جدول ۲{-}۱)}

% SOURCE Synthesis_Research_Gap_Chapter2.txt:4-4 [paragraph]
برای ارائه مقایسه یکپارچه، تحلیل ساختارمند میان ۱۰ مقاله مرجع بررسی‌شده و تبیین شفاف جایگاه روش‌شناختی پژوهش حاضر در ادبیات تحقیق، ماتریس مقایسه‌ای پیشینه پژوهش تدوین گردیده است. این ماتریس، مقالات پیشین را از نظر مدالیته تصاویر پزشکی، الگوریتم‌های اصلی، شاخص‌های کمّی ارزیابی (شامل آنتروپی شانون، \lr{NPCR}، \lr{UACI}، \lr{PSNR} و \lr{SSIM})، زمان اجرا و مهم‌ترین محدودیت‌های عملیاتی دسته‌بندی می‌نماید.\par

% SOURCE Synthesis_Research_Gap_Chapter2.txt:5-5 [spacing]


% SOURCE Synthesis_Research_Gap_Chapter2.txt:6-6 [image]
\SourcePicture{chapter2_taxonomy_table.png}{جدول ۲{-}۱: ماتریس جامع مقایسه‌ای کارهای پیشین و مقایسه با ویژگی‌های چارچوب پیشنهادی}{[جدول ۲{-}۱: ماتریس جامع مقایسه‌ای کارهای پیشین و مقایسه با ویژگی‌های چارچوب پیشنهادی {-} \lr{chapter2\_taxonomy\_table.png}]}{1}

% SOURCE Synthesis_Research_Gap_Chapter2.txt:7-7 [spacing]


% SOURCE Synthesis_Research_Gap_Chapter2.txt:8-8 [heading]
\SourceHeading{2}{۲{-}۴{-}۲. تحلیل نقدگونه و تبیین شکاف پژوهشی \lr{(Research Gap)}}

% SOURCE Synthesis_Research_Gap_Chapter2.txt:9-9 [paragraph]
بررسی جامع ادبیات تحقیق و کالبدشکافی ماتریس مقایسه‌ای فوق، نشان می‌دهد که پژوهش‌های پیشین عمدتاً بر روی یکی از ابعاد مسئله (یا صرفاً رمزنگاری تصویر و یا صرفاً پنهان‌نگاری فراداده) تمرکز داشته‌اند و نتوانسته‌اند یک زنجیره امنیتی یکپارچه، متوازن و آگاه به محدودیت‌های بالینی ارائه نمایند. بر این اساس، سه شکاف پژوهشی اساسی \lr{(Research Gaps)} در ادبیات تحقیق استخراج گردیده‌اند که شالوده منطقی شکل‌گیری اهداف و فرضیات این پایان‌نامه را تشکیل می‌دهند \lr{[1{-}10]}:\par

% SOURCE Synthesis_Research_Gap_Chapter2.txt:10-10 [spacing]


% SOURCE Synthesis_Research_Gap_Chapter2.txt:11-11 [paragraph]
۱. شکاف تولید کلید پویا و پایداری در برابر حملات متن‌واضح برگزیده \lr{(CPA/KPA)}:\par

% SOURCE Synthesis_Research_Gap_Chapter2.txt:12-12 [paragraph]
اکثر سامانه‌های رمزشده آشوبی سنتی از کلیدهای ایستا یا وابسته به پارامترهای ثابت استفاده می‌کنند که سیستم را در برابر حملات متن‌واضح برگزیده \lr{(Chosen{-}Plaintext Attack)} و متن‌واضح معلوم \lr{(Known{-}Plaintext Attack)} بسیار آسیب‌پذیر می‌سازد \lr{[2, 3]}. در پژوهش‌هایی که کلید پویا پیشنهاد گردیده است (مانند استفاده از \lr{EXIF Metadata})، مکانیزم انتقال شرایط اولیه به گیرنده شفاف نبوده یا خودِ فراداده‌های هویتی در طول مسیر ارتباطی بدون محافظت رها شده‌اند \lr{[8]}. علاوه بر این، وابستگی کلید به تمام تصویر باعث می‌شود کوچک‌ترین نویز در پس‌زمینه تصویر واضح، مانع رمزگشایی صحیح گردد.\par

% SOURCE Synthesis_Research_Gap_Chapter2.txt:13-13 [spacing]


% SOURCE Synthesis_Research_Gap_Chapter2.txt:14-14 [paragraph]
۲. شکاف تقابل میان ظرفیت پنهان‌نگاری فراداده و کیفیت تشخیصی بالینی \lr{(PSNR vs ROI)}:\par

% SOURCE Synthesis_Research_Gap_Chapter2.txt:15-15 [paragraph]
روش‌های پنهان‌نگاری موجود (اعم از الگوریتم‌های مبتنی بر یادگیری عمیق مانند \lr{SE{-}Inception} و \lr{Mask{-}RCNN} یا روش‌های کلاسیک \lr{LSB/PVD})، فراداده‌های بیمار را بدون توجه دقیق به ارزش بالینی نواحی تصویر جای‌دهی می‌کنند \lr{[4, 6, 7]}. این امر منجر به ایجاد نویز دیداری، اعوجاج آناتومیک و افت افت کیفیت تشخیصی \lr{(PSNR \SourceGlyph{003C} 40 dB)} در ناحیه حساس تشخیصی \lr{(ROI)} می‌گردد. از سوی دیگر، روش‌هایی که کیفیت بالایی ارائه می‌دهند، فاقد ماژول رمزنگاری آشوبی برای خودِ تصویر تشخیصی هستند \lr{[5]}.\par

% SOURCE Synthesis_Research_Gap_Chapter2.txt:16-16 [spacing]


% SOURCE Synthesis_Research_Gap_Chapter2.txt:17-17 [paragraph]
۳. شکاف عدم یکپارچه‌سازی لایه‌های امنیتی در سامانه‌های اینترنت اشیای پزشکی \lr{(IoMT)}:\par

% SOURCE Synthesis_Research_Gap_Chapter2.txt:18-18 [paragraph]
هیچ‌یک از کارهای پیشین نتوانسته‌اند قطعه‌بندی معنایی ناحیه \lr{ROI} با شبکه \lr{U{-}Net}، رمزنگاری ابرآشوبی ۵ بعدی با بیش از یک توان لیاپانوف مثبت و انتشار \lr{DNA}، و پنهان‌نگاری آگاه به برجستگی بصری \lr{(Saliency{-}aware Steganography)} فراداده‌های \lr{zlib} را در قالب یک معماری یکپارچه و سبک با زمان اجرای زیر ۳ ثانیه برای پلتفرم‌های رایانش لبه (\lr{Edge Computing} در \lr{IoMT}) مجتمع سازند \lr{[1, 2, 6]}.\par

% SOURCE Synthesis_Research_Gap_Chapter2.txt:19-19 [spacing]


% SOURCE Synthesis_Research_Gap_Chapter2.txt:20-20 [paragraph]
چارچوب هیبریدی پیشنهادی در این پایان‌نامه دقیقاً برای پاسخگویی به این شکاف‌های سه‌گانه و اثبات فرضیات پژوهش تدوین گردیده است که جزئیات دقیق ریاضی و الگوریتمی آن در فصل سوم ارائه خواهد شد.\par

% SOURCE Synthesis_Research_Gap_Chapter2.txt:21-21 [spacing]

